\documentclass[11pt]{amsart}
\usepackage[margin=1.1in]{geometry}
\usepackage{amsmath,amssymb,amsthm,booktabs,array,enumitem,xcolor,hyperref}
\hypersetup{colorlinks=true,linkcolor=blue!50!black,citecolor=blue!50!black,urlcolor=blue!50!black}
\newtheorem{theorem}{Theorem}[section]
\newtheorem{proposition}[theorem]{Proposition}
\newtheorem{lemma}[theorem]{Lemma}
\newtheorem{corollary}[theorem]{Corollary}
\newtheorem{conjecture}[theorem]{Conjecture}
\theoremstyle{definition}
\newtheorem{definition}[theorem]{Definition}
\newtheorem{remark}[theorem]{Remark}
\newtheorem{computation}[theorem]{Computation}
\renewcommand{\O}{\mathbb{O}}\newcommand{\Sd}{\mathbb{S}}\newcommand{\Hq}{\mathbb{H}}\newcommand{\Z}{\mathbb{Z}}\newcommand{\Q}{\mathbb{Q}}\newcommand{\R}{\mathbb{R}}\newcommand{\C}{\mathbb{C}}\newcommand{\F}{\mathbb{F}}
\newcommand{\Stab}{\operatorname{Stab}}\newcommand{\Aut}{\operatorname{Aut}}\newcommand{\Der}{\operatorname{Der}}\renewcommand{\Im}{\operatorname{Im}}
\title[Zero divisors of the integral sedenions]{Zero divisors of the integral sedenions\\ and the compact $G_2$\\[4pt] {\normalsize Part IV of the Octonion Program record}}
\author{Christopher Younge}
\address{Philadelphia, PA}
\date{4 October 2026 --- working draft 8, not for circulation}
\begin{document}
\begin{abstract}
We study the zero divisors of the integral sedenions, the Cayley--Dickson double of Coxeter's order of integral octonions, and their interaction with the arithmetic of the compact group $G_2$. Among the $57{,}600$ sedenions whose two octonion halves are units, exactly $7{,}560$ are zero divisors, namely those whose halves are imaginary and orthogonal; they fall into two $G_2(\Z)$-orbits of sizes $6{,}048$ and $1{,}512$, distinguished by the type (Lipschitz or Hurwitz) of the quaternion subalgebra they generate, and the $84{,}672$ ordered pairs $(x,y)$ with $xy=0$ form seven free $G_2(\Z)$-orbits: the integral shadow of Moreno's theorem that the real unit zero-divisor pairs form the Lie group $G_2$. Over $\R$ the zero-divisor manifold is $G_2/\mathrm{Sp}(1)$ and its Hecke algebra $\mathcal H(G_2,\mathrm{Sp}(1))$ is computed in low weights. The conic of a quaternion subalgebra modulo $p$ meets the maximal-parahoric class set of the compact $G_2$ at $p$ in a pattern governed by the octahedral group $S_4$: for $p\in\{3,5,7,11\}$ each conic lies in a single class determined by the type of the zero divisor, and for $p\ge 13$ it splits into $S_4$-orbits, with class stabilisers $192/m$ and $48/m$; this is proved for all $p\ge 5$. Zero-divisor periods of algebraic modular forms on $G_2$ at level~3 are defined; they are the $D^1$-periods of Gan and Savin's exceptional see-saw, and the period along the conic separates Iwahori-new forms from the old ones (through the degeneracy map), which are identified exactly as Gan--Gurevich-type lifts of the level-$3$ newforms of weights $6$ and $10$ with the classical form on $D^1$. Three constructions of a sedenion spectral geometry in the sense of Farnsworth and Boyle are shown to admit no finite Dirac operator beyond a scalar mass, and the Euler product of the norm-counting zeta function is shown to survive exactly up to the sedenions along the Cayley--Dickson ladder. All computations and code are deposited in the program's repository.
\end{abstract}
\maketitle
\tableofcontents

\section{Introduction and summary}
Parts I and II of this record \cite{PartI,PartII} concern class sets and Hecke modules of the compact groups $G_2$ and $F_4$ built from Coxeter's integral octonions $\O_\Z$ (the $E_8$ lattice). The present part goes one rung further on the Cayley--Dickson ladder, to the sedenions $\Sd=\O\oplus\O\ell$, and asks what the first failure of division has to do with the arithmetic developed so far. The answer is more structured than expected.

\begin{enumerate}[leftmargin=2em]
\item[(a)] The integral sedenions $\Sd_\Z=\O_\Z\oplus\O_\Z\ell$ form the lattice $E_8\oplus E_8$. Their zero divisors with unit halves are \emph{exactly} the elements $a+b\ell$ with $a,b$ imaginary orthogonal units (Theorem~\ref{thm:zd}); there are $7{,}560$ of them, each with a $4$-dimensional annihilator, and they form two $G_2(\Z)$-orbits separated by the type of the quaternion subalgebra $\langle 1,a,b,ab\rangle$. The ordered zero-divisor pairs form seven free $G_2(\Z)$-orbits (Theorem~\ref{thm:seven}).
\item[(b)] Over $\R$ the unit zero divisors form $G_2/\mathrm{Sp}(1)$ and Moreno's $G_2$ is the sphere bundle of their annihilators. The Hecke algebra $\mathcal H(G_2,\mathrm{Sp}(1))$, a non-commutative algebra with continuous families of idempotents, has blocks of sizes $1,3,3,6,8,10,6$ on the representations of dimensions $1,7,14,27,64,77,77'$ (Section~\ref{sec:hecke}).
\item[(c)] The conic law (Theorem~\ref{thm:conic}): the isotropic lines of $\Im\Hq_z\otimes\F_p$ form a conic $\cong\mathbb P^1(\F_p)$, and its intersection with the maximal-parahoric class set of the compact $G_2$ at $p$ is governed by the octahedral group $S_4$. The two zero-divisor types index two classes for $p=3,5,7,11$ and the conics split into cube-orbits for $p\ge 13$.
\item[(d)] Zero-divisor periods of algebraic modular forms (Section~\ref{sec:periods}) are the $D^1$-periods of the Gan--Savin see-saw \cite{GanSavin}; the period along the conic vanishes exactly on the Iwahori-new forms of weight $2\omega_1$ at level~3, while their periods on other orbits do not.
\item[(e)] Negatives (Section~\ref{sec:negatives}): in the Farnsworth--Boyle framework of non-associative spectral geometry \cite{FB13,BF14} the sedenions admit, in all constructions tried, no finite Dirac operator beyond a scalar mass; the sedenion analogue of the Albert algebra is not Jordan; and the Euler product of the norm-counting zeta function dies one doubling past the sedenions, where Ramanujan's $\Delta$ first enters.
\end{enumerate}
Everything in (a)--(d) was computed exactly (integer arithmetic, finite groups, or linear algebra at machine precision with independent checks); the conic law is proved for all $p\ge 5$ by a structural reduction to a finite computation (Section~\ref{sec:conic}).

\section{The integral sedenions and their zero divisors}\label{sec:zd}
Let $\O_\Z$ be Coxeter's maximal order of integral octonions, with norm $N$, trace $\mathrm{tr}$ and the $240$ units. The Cayley--Dickson double $\Sd=\O\oplus\O\ell$ has product $(a,b)(c,d)=(ac-\bar d b,\ da+b\bar c)$, conjugation $\overline{(a,b)}=(\bar a,-b)$ and norm $N(a,b)=N(a)+N(b)$; we take $\Sd_\Z=\O_\Z\oplus\O_\Z\ell$. On $200$ random triples the laws behave as classically known: $\Sd$ is flexible and power-associative; it is not commutative, associative or alternative, and $N(xy)/(N(x)N(y))$ ranged over $[0.709,1.316]$.

\begin{theorem}\label{thm:zd}
Let $x=(a,b)\in\Sd_\Z$ with $N(a)=N(b)=1$. Then $x$ is a zero divisor if and only if $a$ and $b$ are both imaginary ($\mathrm{tr}=0$) and orthogonal. There are $7{,}560$ such elements; each has a $4$-dimensional annihilator $\{y: xy=0\}$. No unit of $\Sd_\Z$ (norm~$1$) is a zero divisor.
\end{theorem}
\begin{proof}
Exhaustive computation over all $240^2=57{,}600$ pairs $(a,b)$ of units: the left-multiplication matrix $L_x$ is singular exactly for the $7{,}560$ pairs with $\mathrm{tr}\,a=\mathrm{tr}\,b=0$ and $\langle a,b\rangle=0$ (the other seven combinations of the three conditions never give a zero divisor), and its kernel has dimension $4$ in every case. A unit has one half zero and all $230{,}400$ products of units have norm~$1$.
\end{proof}
The real version (imaginary, orthogonal, equal norm) is due to Moreno and to Khalil--Yiu; the count and the dimension $4$ reproduce Moreno's dimension count $11+3=14=\dim G_2$ for the set of unit zero-divisor pairs \cite{Moreno}.

\begin{theorem}\label{thm:seven}
The group $G_2(\Z)=\Aut(\O_\Z)$ (order $12{,}096$) acts on the $7{,}560$ zero divisors with two orbits, of sizes $6{,}048$ and $1{,}512$, with stabilisers of orders $2$ and $8$. The quaternion subalgebra $\Hq_z=\langle 1,a,b,ab\rangle\cap\O_\Z$ contains $8$ units (Lipschitz type) on the large orbit and $24$ (Hurwitz type) on the small one, and these are exactly the numbers of zero divisors $y$ in the annihilator of $x$ among the $7{,}560$. The $84{,}672=7\cdot 12{,}096$ ordered pairs $(x,y)$ of these zero divisors with $xy=0$ form seven free $G_2(\Z)$-orbits.
\end{theorem}
\begin{proof}
Computation with the $12{,}096$ automorphisms in the Coxeter basis (each verified to be a sedenion automorphism acting diagonally on the two halves); the orbit and stabiliser counts are exact. The map $(a,b)\mapsto(a,-b)$ is an automorphism of $\Sd$; the half-swap preserves the zero-divisor set but is not an automorphism.
\end{proof}
\begin{remark}
Moreno's theorem says the real unit zero-divisor pairs form a space homeomorphic to $G_2$; Theorem~\ref{thm:seven} is its integral shadow: seven copies of $G_2(\Z)$, acted on freely.
\end{remark}

\section{The zero-divisor manifold and its Hecke algebra}\label{sec:hecke}
Over $\R$, a unit zero divisor is a pair $(a,b)$ of orthonormal imaginary octonions. $G_2$ acts transitively on such pairs, so the zero-divisor manifold is $ZD=G_2/K$ with $K$ the stabiliser of one pair.
\begin{proposition}
$K$ is $3$-dimensional with Lie algebra $\mathfrak{su}(2)$; it fixes $\Hq_z$ pointwise and acts on the $4$-dimensional annihilator of $z=(a,b)$ as $\mathrm{Sp}(1)$ acts on $\Hq$ (the commutant of $K$ on the annihilator is $4$-dimensional). Hence $ZD=G_2/\mathrm{Sp}(1)$ is the Stiefel manifold of orthonormal pairs in $\R^7$, of dimension $11$, and Moreno's $G_2$ is the unit-sphere bundle of the annihilators over it.
\end{proposition}
\begin{proof}
Computed from the $14$ derivations of $\O$ in an orthonormal basis: the derivations killing $a$ and $b$ form a $3$-dimensional subalgebra, which kills $\langle 1,a,b,ab\rangle$ and preserves the annihilator, on which its commutant has dimension $4$.
\end{proof}
In Gan--Savin's notation \cite{GanSavin}, $K=D^1$ is the group of norm-one elements of the quaternion subalgebra $D=\Hq_z$, acting by $g\cdot(a,b)=(a,b\bar g)$; it is the pointwise stabiliser of $D$ in $\Aut(\O)$.

The Hecke algebra of the homogeneous space is $\mathcal H(G_2,K)=\bigoplus_\pi M_{m_\pi}(\C)$ with $m_\pi=\dim\pi^K$ (Frobenius reciprocity: $m_\pi$ is the multiplicity of $\pi$ in $L^2(ZD)$). Decomposing $7\otimes 7$, $7\otimes 14$, $\mathrm{Sym}^3 7$ and $\mathrm{Sym}^2 14$ by the Casimir and counting $K$-fixed vectors:
\begin{center}
\begin{tabular}{l@{\quad}ccccccc}\toprule
$\pi$ & $1$ & $7$ & $14$ & $27$ & $64$ & $77$ & $77'$\\
$m_\pi=\dim\pi^K$ & $1$ & $3$ & $3$ & $6$ & $8$ & $10$ & $6$\\\bottomrule
\end{tabular}
\end{center}
The $7$ gives $3$ because $K$ fixes exactly $\Im\Hq_z\subset\Im\O$; the $14$ gives $3$ because the centraliser of $K$ in $\mathfrak g_2$ is the other $\mathfrak{su}(2)$ of $\mathfrak{so}(4)\subset\mathfrak g_2$. Since $(G_2,K)$ is not a Gelfand pair, the idempotents of each block form continuous families (Grassmannians). The representation theory here is classical; the identification of $ZD$ with the sedenion zero divisors is what makes it ours. Every weight in the Path~A table of compact $G_2$ forms (Part~II) has $m_\pi>0$, which is what makes the periods of Section~\ref{sec:periods} possible.

\section{Quaternion subalgebras and the conic law}\label{sec:conic}
\subsection{Quaternion subalgebras of the Coxeter order}
Every zero divisor $z=(a,b)$ determines $\Hq_z=\langle 1,a,b,ab\rangle\cap\O_\Z$, a quaternion order of the definite quaternion algebra it spans. Counting ordered orthogonal imaginary unit pairs per subalgebra ($6\cdot 4=24$ in both cases) gives:
\begin{proposition}\label{prop:count}
$\O_\Z$ contains $63$ quaternion subalgebras of Hurwitz type ($24$ units) and $252$ of Lipschitz type ($8$ units). Their stabilisers in $G_2(\Z)$ have orders $192$ and $48$; the pointwise stabilisers have orders $8$ and $2$ (the stabilisers of the zero divisors of Theorem~\ref{thm:seven}), and in both cases the quotient acting faithfully on $\Im\Hq_z$ is the octahedral group $S_4$.
\end{proposition}
\begin{proof}
$1512/24=63$, $6048/24=252$; $12096/63=192$, $12096/252=48$; the pointwise stabilisers were computed directly, and $192/8=48/2=24=|S_4|$, with $S_4$ the full rotation group of the cube $\{\pm i,\pm j,\pm k\}$, realised in both cases by the computation of Section~\ref{sec:periods} (orbit sizes $6,8,12$ below).
\end{proof}

\begin{proposition}\label{prop:e7}
The imaginary integral units are the $126$ roots of the $E_7$ lattice $\Im\O_\Z$, and the $7{,}560$ zero divisors of Theorem~\ref{thm:zd} are exactly the ordered pairs of orthogonal $E_7$ roots ($60$ roots are orthogonal to any given root, a $D_6$). The zero-divisor mass $7560/12096=5/8$ is thus a Siegel representation number of the binary form $2I_2$ by $E_7$ divided by $|G_2(\Z)|$, and the zero divisors of each norm are governed by Siegel's local--global formula for representations by $E_7$, even though their counting series has no Euler product of its own.
\end{proposition}
\begin{proposition}\label{prop:seven}
For a zero-divisor pair $x=(a,b)$, $y=(c,d)$ with $xy=0$ one has $d=\pm(ac)b$, so the pair is the triple $(a,b,c)$ with $c$ a unit perpendicular to $\Hq_x$. The seven free $G_2(\Z)$-orbits of Theorem~\ref{thm:seven} are separated by the types (Hurwitz $24$ or Lipschitz $8$) of the five quaternion subalgebras $\langle a,b\rangle,\langle c,d\rangle,\langle a,c\rangle,\langle b,c\rangle,\langle ab,c\rangle$: the seven realised type-patterns are $(8,8,8,24,8)$, $(8,8,24,8,8)$, $(8,24,8,8,24)$, $(8,24,24,24,24)$ for Lipschitz $x$ and $(24,8,8,24,8)$, $(24,8,24,8,8)$, $(24,24,8,8,24)$ for Hurwitz $x$, with equal counts within each type as the freeness requires.
\end{proposition}

\subsection{The conic}
Fix an odd prime $p$. The space $\Im\O\otimes\F_p$ is a nondegenerate $7$-dimensional quadratic space; its isotropic lines number $(p^6-1)/(p-1)$ and form the maximal-parahoric level-$p$ structures of the compact $G_2$, so that the class set is $X_p=G_2(\Z)\backslash\{\text{isotropic lines}\}$ (Part~I). For a zero divisor $z$, $\Im\Hq_z\otimes\F_p$ is a nondegenerate ternary space, so its isotropic lines form a conic $C_z\cong\mathbb P^1(\F_p)$ with $p+1$ points, on which $\Stab(\Hq_z)$ acts through $S_4\subset\mathrm{PGL}_2(\F_p)$.

\begin{theorem}[Conic law]\label{thm:conic}
Let $p\ge 5$ and $z$ a zero divisor. The $G_2(\Z)$-classes of the points of $C_z$ are exactly the $S_4$-orbits on $\mathbb P^1(\F_p)$, whose sizes lie in $\{6,8,12,24\}$ (cube faces, vertices, edges, and generic), and the stabiliser in $G_2(\Z)$ of a conic point in an orbit of size $m$ has order $192/m$ (Hurwitz type) or $48/m$ (Lipschitz type). The number of classes a conic meets is the number of $S_4$-orbits on $\mathbb P^1(\F_p)$, which by Burnside's lemma is
\[\frac{1}{24}\Big((p+1)+3f_{-1}+6f_{-2}+8f_3+6f_4\Big),\qquad f_{-1}=2\text{ if }\big(\tfrac{-1}{p}\big)=1,\ f_{-2}=2\text{ if }\big(\tfrac{-2}{p}\big)=1,\ f_3=2\text{ if }p\equiv1\ (3),\ f_4=2\text{ if }p\equiv1\ (4),\]
and $0$ otherwise (the face, edge and vertex rotations of the cube have $2$ or $0$ fixed points on the conic according to these conditions). Equivalently, the conic carries a $6$-point orbit (stabilisers $8$ and $32$) exactly when $p\equiv1\ (4)$, an $8$-point orbit ($6$ and $24$) exactly when $p\equiv1\ (3)$, a $12$-point orbit ($4$ and $16$) exactly when $(\tfrac{-2}{p})=1$, and all remaining points lie in $24$-point orbits ($2$ and $8$); this was verified directly for all $p\le 41$. The orbit structure is thus Eisenstein data of level $24$: the number of classes a conic meets is $\tfrac{p+1}{24}+\tfrac34[p\equiv1\,(4)]+\tfrac12[(\tfrac{-2}{p})=1]+\tfrac23[p\equiv1\,(3)]$, and its Dirichlet series over $p$ is a combination of $\zeta(s-1)$, $\zeta(s)$ and the $L$-functions of the characters of conductors $3$, $4$ and $8$, while the cuspidal content of the zero divisors lies entirely in the periods of Section~\ref{sec:periods}. In particular each conic lies in a single class exactly for $p\in\{5,7,11,23\}$ ($p=23$ is the regular orbit, $24$ points with stabilisers $2$ and $8$), and for $p\ge 29$ never; for $p=3$ the conic ($4$ points) also lies in a single class, with the exceptional stabilisers $432$ and $36$. The count depends only on $p$ modulo $24$.
\end{theorem}
\begin{lemma}\label{lem:unique}
For $p\ge 5$, an isotropic line lying on the conic of a quaternion subalgebra of given type lies on no other conic of that type.
\end{lemma}
\begin{proof}
Let $\Hq_1\ne\Hq_2$ be suborders of the same type, $M_i=\Im\Hq_i$ (rank $3$). Since $\Hq_i=B_i\cap\O_\Z$ for the $\Q$-subalgebra $B_i$ it spans, $B_1\ne B_2$, and two distinct quaternion subalgebras of $\O$ meet in a subalgebra of dimension $\le 2$ (a $3$-dimensional subalgebra of a quaternion algebra does not exist); hence $M_1\cap M_2$ has rank $\le 1$. The conics $C_1,C_2$ modulo $p$ share a point only if $\bar M_1\cap\bar M_2\subset\Im\O\otimes\F_p$ contains a nonzero isotropic vector. Now $\dim_{\F_p}(\bar M_1\cap\bar M_2)=6-\operatorname{rank}_p(M_1+M_2)$, and $\operatorname{rank}_p(M_1+M_2)=\operatorname{rank}(M_1+M_2)$ unless $p$ divides the index of $M_1+M_2$ in its saturation. So either (i) $p$ divides that index, or (ii) $\bar M_1\cap\bar M_2$ is the reduction of $M_1\cap M_2$, which is $0$ or $\Z v$ with $v$ a common imaginary vector, isotropic mod $p$ only if $p\mid N(v)$.

Computation over all $252\cdot 251/2=31{,}626$ Lipschitz pairs and $63\cdot 62/2=1{,}953$ Hurwitz pairs (code \texttt{g2table/lemma43.py}): the saturation indices take the values $1,2,3,4,5,6,8,10$ (Lipschitz) and $1,2,3,4,5,6,7,8$ (Hurwitz), so their prime factors lie in $\{2,3,5,7\}$; the pairs with $\operatorname{rank}(M_1+M_2)=5$ ($6{,}426$ Lipschitz, $1{,}197$ Hurwitz) have common directions of norm $1$ or $2$ (Lipschitz) and $1$ or $3$ (Hurwitz). Hence for $p\ge 11$ case (i) never occurs and case (ii) gives no isotropic vector: the conics are disjoint. For $p=5$ and $p=7$, case (i) occurs for $15+75$ pairs at $p=5$ and $23$ Hurwitz pairs at $p=7$; for each of these the intersection $\bar M_1\cap\bar M_2$ was computed over $\F_p$ and found to contain no isotropic vector, while case (ii) is excluded as before since $p\nmid N(v)\in\{1,2,3\}$. This proves the lemma for all $p\ge 5$. (For $p=3$ it fails: $3\mid N(v)$ for the Hurwitz pairs with a common direction of norm $3$, and the conics do share points, which is why the stabilisers $432,36$ at $p=3$ exceed $192/4$ and $48/4$.)
\end{proof}
\begin{proof}[Proof of Theorem \ref{thm:conic} from Lemma \ref{lem:unique}]
Let $\ell\in C_z$ and $g\in G_2(\Z)$ with $g\ell\in C_z$. Then $g\ell$ lies on the conics of $\Hq_z$ and of $g\Hq_z$, both of the type of $z$; by the lemma $g\Hq_z=\Hq_z$, so $g\in\Stab(\Hq_z)$. Hence the $G_2(\Z)$-orbit of $\ell$ meets $C_z$ exactly in its $\Stab(\Hq_z)$-orbit, which is its $S_4$-orbit since the pointwise stabiliser acts trivially on $C_z$, and the stabiliser of $\ell$ in $G_2(\Z)$ equals its stabiliser in $\Stab(\Hq_z)$, of order $|\Stab(\Hq_z)|/m$. The orbit sizes of $S_4\subset\mathrm{PGL}_2(\F_p)$ on $\mathbb P^1(\F_p)$ are $24/|$point stabiliser$|$ with point stabilisers the cyclic subgroups of orders $4,3,2,1$, giving $6,8,12,24$; $S_4$ is transitive on $\mathbb P^1(\F_p)$ iff $p+1\in\{6,8,12\}$ (the case $4$ being $p=3$, where the lemma fails and the stabilisers are larger).
\end{proof}
\begin{computation}
Independently of the proof, exact orbit computations for $p=3,5,7,11,13,17,19$ (four to forty zero divisors of each type per prime) give the class structure of the conic points shown in Table~\ref{tab:conic}; the stabiliser identity $|{\rm Stab}(\ell)|=|\Stab(\Hq_z)|/m$ holds in every case with $p\ge 5$, as the theorem predicts. The classes touched were also matched against the full class sets $X_5$ ($3$ classes, stabilisers $32,8,6$: the $6$-class is touched by no zero divisor) and $X_7$ ($8$ classes).
\end{computation}
\begin{table}[h]\centering
\begin{tabular}{c c l l}\toprule
$p$ & $p+1$ & Lipschitz: (points, stabiliser) & Hurwitz: (points, stabiliser)\\\midrule
3 & 4 & $(4,36)$ & $(4,432)$\\
5 & 6 & $(6,8)$ & $(6,32)$\\
7 & 8 & $(8,6)$ & $(8,24)$\\
11 & 12 & $(12,4)$ & $(12,16)$\\
13 & 14 & $(8,6)+(6,8)$ & $(8,24)+(6,32)$\\
17 & 18 & $(12,4)+(6,8)$ & $(12,16)+(6,32)$\\
19 & 20 & $(12,4)+(8,6)$ & $(12,16)+(8,24)$\\\bottomrule
\end{tabular}
\caption{Class structure of the conic points of a zero divisor; stabilisers are $48/m$ and $192/m$ for $p\ge 5$.}\label{tab:conic}
\end{table}
\begin{remark}
The incidence counts are now explained: $63(p+1)$ Hurwitz and $252(p+1)$ Lipschitz incidences fall on classes of sizes $12096\,m/192$ and $12096\,m/48$, one conic through each point.
\end{remark}

\section{Zero-divisor periods}\label{sec:periods}
\subsection{Definition}
Let $f$ be an algebraic modular form of weight $\pi$ for the compact $G_2$ at a parahoric level at $p=3$: a function on the level-$3$ structures $x$ (isotropic lines, planes or flags mod~$3$) with $f(x\gamma)=\pi(\gamma)^{-1}f(x)$, in the conventions of \cite{PartII}. For a zero divisor $z$ with continuous stabiliser $K_z$, and $u\in\pi^{K_z}$, put
\[P_{z,u}(f)=\sum_{x\ \text{compatible with }z}\langle f(x),u\rangle,\qquad x\text{ compatible}\iff\text{the line of }x\text{ lies in }\Im\Hq_z\bmod 3,\]
and more generally, for any orbit $O$ of the finite group $\Stab(\Hq_z)$ on the structures, $P_{z,O}(f)=\mathrm{proj}_{\pi^{K_z}}\sum_{x\in O}f(x)$. Every zero divisor has exactly $4$ compatible lines (the conic at $p=3$).

\subsection{Trivial weight}
At the maximal parahoric level $X_3$ has two classes (stabilisers $432,36$) and $T_2$-eigenvalues $126$ (Eisenstein) and $9$ (the lift of $3.6.a.a$). By Theorem~\ref{thm:conic} the four conic points of a Hurwitz zero divisor lie in the $432$-class and those of a Lipschitz one in the $36$-class, so $P_z(f)=4f(\text{class of }z)$: for the cusp form, normalised $f=(-1,1/12)$, the periods are $-4$ and $1/3$.

\subsection{Weights 14 and 27}
Weight $14$ (adjoint): one form at level~$3$ ($T_2=-18$, at the plane and flag levels), whose conic periods vanish for both types. Weight $27$ ($2\omega_1$): at the Iwahori level a $7$-dimensional space with $T_2\in\{40.5\ (\times 2),\ 0\ (\times 3),\ -10.768,\ 6.268\}$, the irrational pair being the roots of $x^2+4.5x-67.5$ and new at the Iwahori level; the others descend to the maximal parahoric levels, where $(T_2,T_5)=(81/2,\,2034/5)$ and $(0,\,11466/25)$ (twice), computed with the calibration $(9,810)$ and $(126,19530)$ in trivial weight.

\begin{computation}\label{comp:periods}
(i) The conic period $P_{z,u}$ vanishes identically on the Iwahori-new forms and is nonzero on all old forms, for both types. (ii) Over all $\Stab(\Hq_z)$-orbits ($51$ orbits for Lipschitz, $21$ for Hurwitz type), the Iwahori-new forms have nonzero $P_{z,O}$ on $41/51$ resp.\ $16/21$ orbits, vanishing exactly on the orbits containing the conic flags (and a few others); the old forms are nonzero on every orbit. (iii) For the Iwahori-new forms the sum of $f$ over the four flags above \emph{each} conic point vanishes separately, before any projection.
\end{computation}
\begin{remark}[what the conic vanishing is]
Item (iii) explains item (i) completely: the sum of $f$ over the planes containing a fixed line is the degeneracy map from the Iwahori level to the maximal-parahoric line level, and Iwahori-new forms lie in its kernel for \emph{every} line. The conic period factors through this degeneracy map, so its vanishing on new forms is the definition of newness, not a property of the zero divisors; the Hurwitz/Lipschitz type and the $K$-projection play no role in it. The content of zero-divisor periods is therefore (a) at the maximal parahoric level, where the conic law (Theorem~\ref{thm:conic}) governs them, and (b) in orbit-periods $P_{z,O}$ that do not factor through degeneracy maps, which is where the new forms are nonzero, in accordance with the global expectation of the next subsection.
\end{remark}
\subsection{The maximal parahoric level $7$ and the selection rule}
At the maximal parahoric level $p=7$ the class set has $8$ classes (stabilisers $2,18,42,6,6,2,6,24$), and the conics of the Hurwitz and Lipschitz zero divisors lie in the $24$-class and in one $6$-class (Theorem~\ref{thm:conic}). Since the $p+1$ conic points form a single orbit of the finite group $\Stab(\Hq_z)$ when $p\le 11$, the conic sum $\sum_{\ell\in C_z}f(\ell)$ is the projection of $f(\ell_0)$ onto the $\Stab(\Hq_z)$-fixed vectors of $V_\pi$, which gives a selection rule: $\dim V_\pi^{\Stab(\Hq_z)}$ is $1$ for the trivial weight, $0$ for the $7$ and the $14$, and $1$ (Hurwitz) or $2$ (Lipschitz) for the $27$. The rule for all weights up to $77$, computed from the characters of $\Stab(\Hq_z)$ on the eigenvalues of its elements on the $7$, is
\begin{center}\begin{tabular}{l@{\quad}ccccccc}\toprule
weight & $1$ & $7$ & $14$ & $27$ & $64$ & $77$ & $77'$\\
$\dim V_\pi^{\Stab(\Hq_z)}$, Lipschitz ($|\Stab|=48$) & $1$ & $0$ & $0$ & $2$ & $1$ & $1$ & $3$\\
$\dim V_\pi^{\Stab(\Hq_z)}$, Hurwitz ($|\Stab|=192$) & $1$ & $0$ & $0$ & $1$ & $0$ & $0$ & $2$\\\bottomrule\end{tabular}\end{center}
so from weight $64$ on the two types differ: in weights $64$ and $77$ every form is dark on Hurwitz conics and (generically) lit on Lipschitz ones. This was verified in weight $64$ at level $7$ ($97$ forms, the $64$ cut out of $7\otimes14$ by the character projector over $G_2(\Z)$): all $97$ conic sums vanish for the Hurwitz type and none for the Lipschitz type. Accordingly all conic sums vanish identically in weights $7$ and $14$ for every form, and in weight $27$ ($50$ forms at this level) the periods are nonzero for every form and both types; in trivial weight ($7$ tempered cusp forms) likewise. So zero-divisor periods do not vanish on tempered forms: they are generic nonzero functionals, in accordance with the nonvanishing predicted by the see-saw, and their content lies in the conic law and the selection rule rather than in any vanishing.

The same computation extends Theorem~\ref{thm:lifts} to level $7$: the five newforms of $S_{10}(\Gamma_0(7))$ appear in weight $2\omega_1$ at the maximal parahoric level with exactly the predicted $(T_2,T_3)$, namely $(58.351,116.103)$, $(37.021,188.629)$, $(35.169,28.970)$ (multiplicity $2$), $(14.331,164.808)$ (multiplicity $2$) and $(1.378,80.601)$; the first two violate the tempered bound, the others pass it. In trivial weight the level-$7$ lifts of weight-$6$ newforms have no fixed vector at this parahoric and live at the plane parahoric.

\subsection{Interpretation}
In \cite{GanSavin} the pair $D^1\subset\Aut(\O)$ sits in a see-saw inside the adjoint $E_7$ of the Albert algebra, with $\Aut(\O)\times\mathrm{PGSp}_6$ and $D^1\times G_D$ ($G_D$ of type $D_6$) as dual pairs; the theta lift of the trivial representation of $D^1$ is a residue of a Siegel Eisenstein series on $G_D$, and by Pollack's integral and Chenevier's identity $L(s,\pi,\mathrm{Spin})=\zeta(s)L(s,\pi',\mathrm{Std})$ for $G_2$-parameters, the $D^1$-period of a form on the compact $G_2$ is governed by $L(1,\pi',\mathrm{Std})\ne 0$. Our periods are $D^1$-periods with specific local data at $3$. Computation~\ref{comp:periods}(ii) is in accordance: the genuine (Iwahori-new) forms have nonzero $D^1$-periods on orbits that do not factor through degeneracy maps; the vanishing along the conic is the degeneracy map itself (Remark above). The old forms are identified. Writing $\mathrm{SO}(4)=(\mathrm{SL}_2^{\rm long}\times\mathrm{SL}_2^{\rm short})/\mu_2\subset G_2$ for the stabiliser of a quaternion subalgebra, with $7=(2,2)\oplus(1,3)$ and $\mathrm{SL}_2^{\rm long}=D^1$, consider the Arthur-type parameter $\psi_f=\varphi_f\boxtimes\nu$ with a classical newform $f$ of weight $k$ on the long-root $\mathrm{SL}_2$ and the Arthur $\mathrm{SL}_2$ on the short root. Since the Hecke eigenvalue of an algebraic modular form is the Satake transform of the unitary local representation (the weight enters only at infinity), $T_p=p^3\chi_7(s_p)-1$ in every weight, and here
\[\chi_7(s_p)=\Big(p+1+\tfrac1p\Big)+\big(p^{1/2}+p^{-1/2}\big)\frac{a_p(f)}{p^{(k-1)/2}},\qquad T_p=\frac{p^3(p^2+p+1)}{p}-1+\frac{(p+1)\,a_p(f)\,p^3}{p^{k/2}}.\]
\begin{theorem}\label{thm:lifts}
At level $3$ the trivial-weight cusp form $(T_2,T_5,T_7)=(9,810,2472)$ is $\psi_{3.6.a.a}$ ($a_2,a_5,a_7=-6,6,-40$), and the weight-$2\omega_1$ old forms are $\psi_{3.10.a.a}$ ($a_2,a_5=18,-1530$; $T_2=81/2$, $T_5=2034/5$) and $\psi_{3.10.a.b}$ ($a_2,a_5=-36,-1314$; $T_2=0$, $T_5=11466/25$, multiplicity $2$ at the maximal parahoric level), with the weight rule $k=6+2a$ for $G_2$-weight $a\omega_1$. All six (resp.\ four) Hecke eigenvalues agree exactly; the prime $7$ was not used in finding the dictionary.
\end{theorem}
These are the Gan--Gurevich non-tempered forms \cite{GG06} seen on the compact $G_2$: the classical form lives on $D^1$, which is why their $D^1$-periods are large, and the two weight-$10$ forms have opposite Atkin--Lehner signs at $3$ ($a_3=\pm 81$), a natural source of the multiplicity difference (to be confirmed against Arthur's multiplicity formula). No level-one forms exist in weights $14$, $27$, $77'$ ($G_2(\Z)$-invariants vanish). The Iwahori-new irrational pair is not of this type and is a genuine (tempered) form.

\section{Negatives and boundaries}\label{sec:negatives}
\subsection{Sedenion spectral geometries}
Farnsworth and Boyle \cite{FB13,BF14} extend Connes' spectral triples to non-associative algebras: the Hilbert space must be a bimodule in the algebra's own class (every identity surviving the replacement of one algebra element by a vector), the order conditions become vanishing associators in $B=\Omega A\oplus H$, and the gauge group is $\Aut(A)$. For the sedenions acting on themselves: the flexible order-zero law holds and the alternative ones fail; the derivation algebra is $14$-dimensional and inner (inside $\mathfrak{so}(16)\oplus\R$); the fermions decompose as $2(1+7)$ under $G_2$ with the $S_3$ of $\Aut(\Sd)=G_2\times S_3$ as a family symmetry; the commutant of the left action is trivial, forcing $\gamma=\pm 1$ and $D_F=0$. On two copies $\Sd\oplus\Sd$ with $\gamma=\mathrm{diag}(1,-1)$, the flexible order-one condition forces $D_F$ to be a multiple of the identity for either real structure; with the opposite bimodule on the second copy, no $D_F$ survives; on the complex sedenions $\C\otimes\Sd$ with real-linear $D_F$, only complex scalars survive. The sedenion analogue $H_3(\Sd)$ of the Albert algebra violates the Jordan identity (relative size $0.96$; $H_3(\O)$ satisfies it to $10^{-15}$), so the Jordan-geometry route \cite{BF14} excludes the sedenions and admits the Albert algebra. In every construction the zero divisors do not enter the action.
\subsection{Euler products along the ladder}
The norm-counting Dirichlet series $\sum r(n)n^{-s}$ has an Euler product iff $r$ is multiplicative. It does for $\Z$, $\Z[i]$, the Hurwitz quaternions ($D_4$), the Coxeter octonions ($E_8$, $r=240\sigma_3$) and the integral sedenions ($E_8\oplus E_8$, $r=480\sigma_7$) --- the last only because $M_8(\mathrm{SL}_2(\Z))$ has no cusp forms --- and fails for the integral $32$-ions ($E_8^4$), since $E_4^4=E_{16}+\tfrac{3456000}{3617}\,E_4\Delta$. The counting zeta function of the $32$-ions, $\tfrac{16320}{3617}\zeta(s)\zeta(s-15)+c\,L(E_4\Delta,s)$, has two zeros on its critical line $\Re s=8$ below $t\approx 15$ and none up to $t=100$, where the Eisenstein part dominates by $10^5$--$10^6$; the Leech lattice's $\zeta(s)\zeta(s-11)-L(\Delta,s)$ has no zeros at all on $\Re s=6$ to $t=120$, its $57$ zeros to height $198$ lying near $\Re s\in[11.1,12.5]$ with spacings more rigid than GUE. The product over primes is lost exactly where $\Delta$ enters.

\section{Open problems}
\begin{enumerate}[leftmargin=2em]
\item A conceptual proof of Lemma~\ref{lem:unique} replacing the finite computation at $p=5,7$ (the primes $5$ and $7$ divide some saturation indices without producing isotropic intersections).
\item The multiplicities of the Arthur-type lifts of Theorem~\ref{thm:lifts} against Arthur's multiplicity formula ($a_3=\pm 81$), and the identification of the Iwahori-new tempered pair through its exceptional theta lift to genus-$3$ Siegel forms of level~$3$.
\item The $32$-ions; the Okubo integral orders of \cite{Corr26} and their idempotent geometry against the conic law.
\item Part III: the $J_E$ half of the $F_4$ class set at level~$5$ ($88\%$ of the mass found; exact deficits $720799/96768$ and $3$).
\end{enumerate}

\appendix
\section{Code and data}
Repository \texttt{chrisyounge-create/Cayley-Agent-Repository}: \texttt{sedenion.py} (doubling, laws), \texttt{sedenion\_geometry.py}, \texttt{sedenion\_step2.py}, \texttt{sedenion\_step3.py}, \texttt{sedenion\_step3\_complex.py} (spectral geometries), \texttt{zd\_hecke.py} (stabiliser, annihilator, multiplicities), \texttt{g2table/zd\_periods.py}, \texttt{g2table/zd\_periods\_vec.py} (periods), \texttt{g2table/lemma43.py} (proof of Lemma~\ref{lem:unique}), \texttt{g2table/zd\_periods\_level7.py} (level-$7$ periods and lifts), \texttt{beurling.py} (prime-correlation model). Dropbox notes \texttt{RH/SEDENIONS.md} through \texttt{RH/NEXT\_RUN\_ITEMS\_1\_TO\_4.md} contain the session records.

\begin{thebibliography}{99}
\bibitem{PartI} C. Younge, \emph{Class sets and Eisenstein congruences on compact exceptional groups}, Part I of this record (2026).
\bibitem{PartII} C. Younge, \emph{The Hecke module of the compact $F_4$ at level~3}, Part II of this record (2026).
\bibitem{Moreno} G. Moreno, The zero divisors of the Cayley--Dickson algebras over the real numbers, Bol. Soc. Mat. Mexicana 4 (1998).
\bibitem{Brown} R. B. Brown, On generalized Cayley--Dickson algebras, Pacific J. Math. 20 (1967).
\bibitem{GG06} W. T. Gan, N. Gurevich, Non-tempered Arthur packets of $G_2$: liftings from $\widetilde{\mathrm{SL}}_2$, Amer. J. Math. 128 (2006).
\bibitem{GanSavin} W. T. Gan, G. Savin, An exceptional Siegel--Weil formula and poles of the spin $L$-function of $\mathrm{PGSp}_6$.
\bibitem{Chen} G. Chenevier, Subgroups of $\mathrm{Spin}(7)$ or $\mathrm{SO}(7)$ with each element conjugate to some element of $G_2$, and applications to automorphic forms, Doc. Math. (2019).
\bibitem{FB13} S. Farnsworth, L. Boyle, Non-associative geometry and the spectral action principle, JHEP 07 (2015) 023, arXiv:1303.1782.
\bibitem{BF14} L. Boyle, S. Farnsworth, Non-commutative geometry, non-associative geometry and the standard model of particle physics, New J. Phys. 16 (2014) 123027.
\bibitem{GG} N. Gillard, N. Gresnigt, Three fermion generations with two unbroken gauge symmetries from the complex sedenions, Eur. Phys. J. C 79 (2019) 446.
\bibitem{Furey} C. Furey, Three generations, two unbroken gauge symmetries, and one eight-dimensional algebra, Phys. Lett. B 785 (2018).
\bibitem{Corr26} D. Corradetti, A. Marrani, Integral orders for the Okubo algebra and idempotent geometries of the $E_8$ lattice, arXiv:2608.27762 (2026).
\bibitem{ConwaySmith} J. H. Conway, D. A. Smith, \emph{On Quaternions and Octonions}, A K Peters (2003).
\end{thebibliography}
\end{document}
